\documentclass[a4paper,12pt]{ctexart}
\usepackage{xcolor}
\usepackage[top=1.8cm, bottom=1.8cm, left=1.8cm, right=1.8cm]{geometry}
\usepackage{array}
\linespread{1.35}

\title{\textcolor{blue!70!black}{\Huge\textbf{从Python到C++}}}
\author{}
\date{}

\begin{document}
\maketitle
\thispagestyle{empty}

\section*{\textcolor{blue!70!black}{说明}}

C++ 是一种更接近计算机底层的高级编程语言，也是信息学竞赛（NOI）的指定语言。从 Python 转到 C++ 需要理解几个关键差异：静态类型 vs 动态类型、手动内存管理、编译 vs 解释。

\textbf{主要内容（初拟）：}
\begin{enumerate}
  \item C++ 与 Python 的核心差异总览
  \item 基本语法对照：变量声明、输入输出
  \item 数据类型与强制类型转换
  \item 控制结构（\texttt{if}、\texttt{for}、\texttt{while}）与 Python 的对照
  \item 数组与指针的初步理解
  \item 函数与函数重载
  \item 面向对象初步：类与对象
  \item 编译与运行：从源代码到可执行文件
  \item 小项目：将 Python 项目改写为 C++
\end{enumerate}

\textbf{典型对照示例：}
\begin{itemize}
  \item Python: \texttt{x = 5} $\leftrightarrow$ C++: \texttt{int x = 5;}
  \item Python: \texttt{if x > 0:} $\leftrightarrow$ C++: \texttt{if (x > 0) \{}
  \item Python: \texttt{print("Hello")} $\leftrightarrow$ C++: \texttt{std::cout << "Hello";}
\end{itemize}

\end{document}
